\part{Structuri de date avansate}

În această parte vom studia structuri de date care, în anumite situații, sînt mai puternice decît arborii de intervale, arborii indexați binar sau structurile de date echilibrate din STL (\ccode{set} sau \ccode{map}). Aceste structuri sînt:

\begin{itemize}
  \item \textit{policy-based data structures} (PBDS);
  \item \textit{skip lists};
  \item treapuri.
\end{itemize}

Ne dorim structuri care stochează chei (ca un \ccode{set}) sau perechi cheie-valoare (ca un \ccode{map}) și care pot răspunde unor nevoi ca:

\begin{enumerate}
  \item \ccode{kth_element(k)}: care este a $k$-a cheie din structură, în ordine crescătoare?
  \item \ccode{order_of(x)}: a cîta cheie este/ar fi $x$, în ordine crescătoare?
  \item Menținerea de chei implicite.
  \item Suport pentru chei mari (neindexabile).
  \item Suport pentru operații complexe, cum ar fi răsturnarea unei porțiuni din structură.
\end{enumerate}

Operațiile 1 și 2 sînt cunoscute colectiv ca \textbf{statistici de ordine} (engl. \textit{order statistics}). Pentru a le implementa, vom introduce conceptul de \textbf{structuri de date augmentate}. În esență, aceste structuri de date rețin, pentru fiecare nod sau pointer, și numărul de noduri acoperite sau subîntinse de acel nod sau pointer.

Cheile implicite se referă la posibilitatea de a menține numerele de ordine ale tuturor elementelor. Ne dorim o structură asemănătoare unui vector (atribuire / citire pe poziție) sau unui arbore de intervale (sume / maxime pe interval), dar care să admită și inserări și ștergeri. Putem vedea aceasta ca pe o structură în care fiecare element are o cheie egală chiar cu indicele elementului. Doar că nu putem stoca \textbf{explicit} acești indici, deoarece la o inserare sau ștergere trebuie să renumerotăm $\bigoh(n)$ indici. De aceea vom stoca \textbf{implicit} cheile într-o manieră care să ne permită să le calculăm eficient la nevoie.

Iată un rezumat al posibilităților oferite de cîteva structuri deja studiate și de cele pe care urmează să le studiem.

\begin{table}[!htbp]
  \centering
  \begin{tabular}{lccccc}
    \textbf{operația} & \textbf{AIB, AINT} & \textbf{set/map STL} & \textbf{PBDS} & \textbf{treap} & \textbf{skip list}\\
    \hline
    inserare & $\bigoh(\log n)$ & $\bigoh(\log n)$ & $\bigoh(\log n)$ & $\bigoh(\log n)$ & $\bigoh(\log n)$\\
    ștergere & $\bigoh(\log n)$ & $\bigoh(\log n)$ & $\bigoh(\log n)$ & $\bigoh(\log n)$ & $\bigoh(\log n)$\\
    căutarea unei chei & $\bigoh(\log n)$ & $\bigoh(\log n)$ & $\bigoh(\log n)$ & $\bigoh(\log n)$ & $\bigoh(\log n)$\\
    minim / maxim & $\bigoh(\log n)$ & $\bigoh(\log n)$ & $\bigoh(\log n)$ & $\bigoh(\log n)$ & $\bigoh(\log n)$\\
    \ccode{kth_element(k)} & $\bigoh(\log n)$ & \xmark & $\bigoh(\log n)$ & $\bigoh(\log n)$ & $\bigoh(\log n)$\\
    \ccode{order_of(x)} & $\bigoh(\log n)$ & \xmark & $\bigoh(\log n)$ & $\bigoh(\log n)$ & $\bigoh(\log n)$\\
    chei implicite & \xmark & \xmark & \xmark & \cmark & \cmark\\
    chei mari & \xmark & \cmark & \cmark & \cmark & \cmark\\
    \hline
  \end{tabular}

  \caption{Structuri de date și posibilitățile lor.}
\end{table}

În experiența mea, PBDS sînt cele mai rapide. Treapurile sînt un pic mai lente, iar \textit{skip lists} cam de 2-3 ori mai lente. Pe de altă parte, consider că \textit{skip lists} sînt mai ușor de implementat decît treap-urile. \textit{YMMV}.

\import{./advanced-ds}{pbds.tex}
\import{./advanced-ds}{skip-lists.tex}
\import{./advanced-ds}{treaps.tex}
